\chapter*{使用说明与星级判定依据}
\addcontentsline{toc}{chapter}{使用说明与星级判定依据}
\markboth{使用说明与星级判定依据}{}

\section*{一、本册收录了哪些题库}

\begin{center}
\begin{tabularx}{\textwidth}{@{}clX@{}}
\toprule
编号 & 来源 & 说明 \\
\midrule
B1 & 《流体力学与流体机械》试题库（四） & 有文字层，选择题 15 题（\textbf{题干括号内直接印有答案}）、问答题 2 题、计算题 4 题 \\
B2 & 西华选择题题库（精简版） & 扫描件，赵琴教材配套选择题库，按章编排共 \textbf{74 题}：绪论 8、流体静力学 10、流体动力学基础 21、管嘴孔口管路水力计算 20、相似理论与量纲分析 10、流体动力学（无旋流动）5；末页是\textbf{手写答案页} \\
B3 & 补充题库（\texttt{17.pdf}） & 百度文库网页截图，实为《流体力学与流体机械》试卷 3，\textbf{只有客观题}：单选 30 题、多选 10 题、判断 20 题，共 60 题 100 分，\textbf{无计算题} \\
\bottomrule
\end{tabularx}
\end{center}

\noindent\textbf{注意：}B1 中含有\textbf{泵与风机（流体机械）}题目。这部分属于教材第 11--14 章，
\textbf{不在 818 初试范围}（初试只考第 1--8 章），本册对这类题已逐题标注"（流体机械，复试范围）"。

\section*{二、星级怎么定的（不是主观印象）}

本册每一道题都标注重要度星级，星级\textbf{只由文件证据决定}：

\begin{center}
\begin{tabularx}{\textwidth}{@{}clX@{}}
\toprule
星级 & 含义 & 判定条件（满足其一） \\
\midrule
\xstarfull{5} & 必做、必会 & 该考点在历年真题中作为\textbf{计算题/大题}出现 $\ge 3$ 次；或考情分析明确写"考试计算题重点" \\
\xstarfull{4} & 高频 & 在历年真题中出现 2--3 次；或考情分析写"重点掌握" \\
\xstarfull{3} & 常考 & 在历年真题中出现 1 次；或考点分析写"掌握基本概念与公式套用" \\
\xstarfull{2} & 了解 & 真题中未出现，但属 818 考纲章节核心内容、且题库中反复出现 \\
\xstarfull{1} & 边缘 & 仅复试范围（气体动力学、流体机械）或纯记忆性补充 \\
\bottomrule
\end{tabularx}
\end{center}

\noindent 每条星级的题干后都附〔依据：$\cdots$〕，指明是\textbf{哪一年的第几题}或\textbf{哪份分析文件的哪一页}。
若一题在真题中原样出现过，会直接写"真题原题：20XX 年第 X 题"。

\section*{三、使用建议}

\begin{enumerate}
  \item \textbf{先做 \xstarfull{5} 与 \xstarfull{4}}：这些题与真题重合度最高，性价比最高；
  \item 选择题要\textbf{把四个选项都弄懂}：818 的选择题考"概念辨析"，
        错误选项往往是另一个考点的正确答案，一个选项就是一次复习；
  \item 计算题按《教材精讲》给出的"解题套路"（选基准面 $\to$ 选断面 $\to$ 定压强基准 $\to$ 列式 $\to$ 代入）
        逐步书写，不要跳步；
  \item 答案统一置于书末，做完再看。
\end{enumerate}
